

% =================================
% Core mathematics
% =================================

\usepackage[utf8]{inputenc}

\usepackage[mathjax,HomeHTMLFilename=index]{lwarp}
\CSSFilename{assets/textbook.css}
% Preserve the original Computer Modern typography in PDF output.
% Lwarp uses T1 encoding for HTML generation.
\ifbool{warpingHTML}{}{%
    \renewcommand{\encodingdefault}{OT1}%
}

\newenvironment{formalderivation}
    {\begin{BlockClass}{formalderivation}}
    {\end{BlockClass}}


% Illustrations, optionally identified for inline links.
\newenvironment{sideillustration}[1][]
  {%
    \begin{BlockClass}{sideillustration}%
    \ifstrempty{#1}{}{\hypertarget{side:#1}{}}%
  }
  {%
    \end{BlockClass}%
  }

\newcommand{\sideLink}[2]
  {\hyperlink{side:#1}{#2}}
    
\newcommand{\bookvspace}[1]{%
  \vspace{#1}%
  \warpHTMLonly{%
    \begin{BlockClass}[height: #1;]{bookvspace}%
    \end{BlockClass}%
  }%
}


\usepackage{amsmath, amssymb, amsthm}
\usepackage{thmtools}

% TeX Live 2026 compatibility fix
\makeatletter
\@ifundefined{newcounteralias}{}{%
  \renewcommand\thmt@autorefsetup{%
    \@xa\def\csname\thmt@envname autorefname\@xa\endcsname
    \@xa{\thmt@thmname}%
  }%
}
\makeatother

\allowdisplaybreaks
\usepackage{mathtools}

% ============================================================
% Page layout and general formatting
% ============================================================

\usepackage[margin=1.25in]{geometry}
\usepackage{parskip}
\usepackage{setspace}
\usepackage{titlesec}
\usepackage{changepage}

\onehalfspacing
\pagenumbering{arabic}
\everymath{\displaystyle}

\setcounter{tocdepth}{5}


% Labels for the navigation links.
\renewcommand{\linkpreviousname}{Previous section}
\renewcommand{\linknextname}{Next section}

% Home page: only a Next link.
\HTMLFirstPageBottom{%
  \begin{BlockClass}{book-pagination}
    \LinkNext
  \end{BlockClass}%
}

% Subsequent pages: Previous and Next where applicable.
\HTMLPageBottom{%
  \begin{BlockClass}{book-pagination}
    \ifnum\value{LWR@htmlseqfilenumber}>1
      \LinkPrevious
    \fi
    \LinkNext
  \end{BlockClass}%
}


% ============================================================
% Colors and small typographical utilities
% ============================================================

\usepackage[dvipsnames]{xcolor}
\usepackage{bm}
\usepackage{nicefrac}
\usepackage{verbatim}
\usepackage{etoolbox}


\usepackage{needspace}

\newcommand{\nbgpara}[2]{%
    \par
    \needspace{3\baselineskip}%
    \addvspace{1.15\baselineskip}%
    \noindent
    \textbf{\S#1}\quad
    \textbf{#2.}\enspace
    \ignorespaces
}

\newcommand{\nbgparar}[1]{%
    \par
    \needspace{3\baselineskip}%
    \addvspace{1.15\baselineskip}%
    \noindent
    \textbf{#1.}\enspace
    \ignorespaces
}

\definecolor{bostonuniversityred}{rgb}{0.8, 0.0, 0.0}

% Long dash.
\newcommand{\srule}{\text{---}}

% Operator names.

\newcommand{\Set}{\operatorname{Set}}
\newcommand{\PC}{\operatorname{PC}}


\begin{warpMathJax}
\CustomizeMathJax{\newcommand{\Set}{\operatorname{Set}}}
\CustomizeMathJax{\newcommand{\PC}{\operatorname{PC}}}



\CustomizeMathJax{%
    \renewcommand{\RightLabel}[1]{%
        \RL{\scriptsize #1}%
    }%
}

\end{warpMathJax}



\newcommand{\Dep}{\operatorname{Dep}}
\newcommand{\Real}{\operatorname{Real}}
\newcommand{\Trans}{\operatorname{Trans}}
% Relations, functions, and basic universes
\DeclareMathOperator{\Dom}{Dom}
\DeclareMathOperator{\Ran}{Ran}
\DeclareMathOperator{\Rel}{Rel}
\DeclareMathOperator{\Fun}{Fun}
\DeclareMathOperator{\Inj}{Inj}
\DeclareMathOperator{\Swell}{Swell}
\DeclareMathOperator{\Basic}{Basic}
\DeclareMathOperator{\FV}{FV}

\newcommand{\defequiv}{\mathrel{\equiv_{\mathrm{df}}}}

% ============================================================
% Section formatting
% ============================================================

\titleformat*{\subsubsection}{\large\bfseries}
\titleformat*{\paragraph}{\color{NavyBlue}\bfseries}

\titleformat{\section}
  {\normalfont\LARGE\bfseries}
  {\thesection}
  {1em}
  {}



\setcounter{tocdepth}{3}
\setcounter{FileDepth}{2}
\setcounter{SideTOCDepth}{3}


\NewDocumentEnvironment{proseSpacing}{O{1.2em} O{1.6}}
  {%
    \begin{BlockClass}[--prose-gap: #1; --prose-leading: #2;]{prose-spacing}%
  }
  {%
    \end{BlockClass}%
  }


% ============================================================
% Diagrams and proof trees
% ============================================================

\usepackage{tikz}
\usetikzlibrary{arrows.meta, positioning, calc}

\usepackage{bussproofs}

% Lwarp compatibility: restore explicit glue for inference lines.
% Inside theorem environments, Lwarp may suppress \hfill,
% which bussproofs requires when constructing inference lines.
\ifbool{warpingHTML}{%
    \patchcmd{\buildSingleScore}
        {\hfill}
        {\hskip 0pt plus 1fill}
        {}
        {\PackageError{bussproofs-fix}
            {Could not patch inference-line construction}
            {Check the bussproofs package implementation}}%
}{}

\let\oldRightLabel\RightLabel
\renewcommand{\RightLabel}[1]{\oldRightLabel{\scriptsize #1}}

\newcommand{\pvspace}[1]{%
    \def\proofSkipAmount{\vskip #1}%
}

\newcommand{\pvreset}{%
    \def\proofSkipAmount{\vskip .8ex plus .8ex minus .4ex}%
}

% ============================================================
% Hyperlinks
% ============================================================

\usepackage{hyperref}
\usepackage{bookmark}

\hypersetup{
  hidelinks
}



% ============================================================
% Shared theorem-style settings
% ============================================================

\newlength{\theoreminset}
\setlength{\theoreminset}{1.5em}

\newlength{\theoremabovespace}
\setlength{\theoremabovespace}{0.2cm}

\newlength{\theorembelowspace}
\setlength{\theorembelowspace}{0.2cm}

% ============================================================
% Theorem-like environments
%
% Shared numbering, reset at each section:
% Definition 1.1, Theorem 1.2, Lemma 1.3, etc.
% ============================================================

\declaretheoremstyle[
  spaceabove=\theoremabovespace,
  spacebelow=\theorembelowspace,
  headfont=\bfseries,
  bodyfont={\normalfont\leftskip=\theoreminset\rightskip=\theoreminset},
  headpunct={},
  headformat={\NAME\ \NUMBER.\NOTE},
  postheadspace=\newline
]{defstyle}


% Additional theorem styles for HTML color coding.
% All share the exact typography and spacing of defstyle.

\newcommand{\DeclareNBGStatementStyle}[1]{%
  \declaretheoremstyle[
    spaceabove=\theoremabovespace,
    spacebelow=\theorembelowspace,
    headfont=\bfseries,
    bodyfont={\normalfont\leftskip=\theoreminset\rightskip=\theoreminset},
    headpunct={},
    headformat={\NAME\ \NUMBER.\NOTE},
    postheadspace=\newline
  ]{#1}%
}

\DeclareNBGStatementStyle{resultstyle}
\DeclareNBGStatementStyle{metastyle}
\DeclareNBGStatementStyle{principlestyle}
\DeclareNBGStatementStyle{neutralstyle}


\declaretheorem[
  style=defstyle,
  name=Theorem,
  numberwithin=section
]{theorem}

\declaretheorem[
  style=metastyle,
  name=Metatheorem,
  sibling=theorem
]{metatheorem}

\declaretheorem[
  style=defstyle,
  name=Definition,
  sibling=theorem
]{definition}

\declaretheorem[
  style=resultstyle,
  name=Proposition,
  sibling=theorem
]{proposition}

\declaretheorem[
  style=resultstyle,
  name=Lemma,
  sibling=theorem
]{lemma}

\declaretheorem[
  style=resultstyle,
  name=Corollary,
  sibling=theorem
]{corollary}

\declaretheorem[
  style=neutralstyle,
  name=Example,
  sibling=theorem
]{example}

\declaretheorem[
  style=neutralstyle,
  name=Remark,
  sibling=theorem
]{remark}

\declaretheorem[
  style=principlestyle,
  name=Principle,
  sibling=theorem
]{principle}

\declaretheorem[
  style=neutralstyle,
  name=Exercise,
  sibling=theorem
]{exercise}


\declaretheorem[
  style=metastyle,
  name={Derived Rule},
  sibling=theorem
]{derivedrule}


% ============================================================
% Axioms
%
% Numbered independently as A_1, A_2, ..., A_8.
% Example heading: A_1   Axiom of Extensionality
% ============================================================

\declaretheoremstyle[
  spaceabove=\theoremabovespace,
  spacebelow=\theorembelowspace,
  headfont=\bfseries,
  bodyfont={\normalfont\leftskip=\theoreminset\rightskip=\theoreminset},
  headpunct={},
  headformat={\NUMBER\quad\NAME\ \NOTE},
  notebraces={}{} ,
  postheadspace=\newline
]{axiomstyle}

\declaretheorem[
  style=axiomstyle,
  name={Axiom of}
]{axiom}

\renewcommand{\theaxiom}{\ensuremath{A_{\arabic{axiom}}}}

% ============================================================
% Special one-off Definition 0 environment
%
% This does not affect the ordinary theorem/definition counter.
% ============================================================

\declaretheoremstyle[
  spaceabove=\theoremabovespace,
  spacebelow=\theorembelowspace,
  headfont=\bfseries,
  bodyfont={\normalfont\leftskip=\theoreminset\rightskip=\theoreminset},
  headpunct={},
  headformat={\NAME.},
  postheadspace=\newline
]{zerodefstyle}

\declaretheorem[
  style=zerodefstyle,
  name={Definition 0 (Informal)},
  numbered=no
]{definitionzero}
